\section{Conclusions and Recommendations}
\label{sec:conclusions}

Road-network car access to resolutive health facilities in Tumbes, Amazonas and Cusco is markedly unequal: a population-weighted mean of \MeanAccessTotal{} minutes overall conceals a range from \MeanAccessTumbes{} minutes in Tumbes to over three times that in parts of rural Cusco and Amazonas, with a Gini coefficient of \GiniCalibrated{} among the population whose access time could be estimated. A substantial share of the frame population (\CovNoTimeTotal\%) currently has no estimable route at all, concentrated in the same districts that also show the longest measured times -- meaning the true picture in those districts is very plausibly worse, not better, than the routed-only figures suggest.

For planning purposes, three points follow directly from the results:
\begin{enumerate}
  \item \textbf{Prioritize by population-weighted gap, not settlement count.} The critical-gap ranking (Table~\ref{tab:worst}) identifies districts such as Camantí, Marcapata and Kosñipata as the largest \emph{measured} gaps; any prioritization exercise should also treat the districts with insufficient data (Section~\ref{sec:gaps}) as a distinct, high-priority data-collection target, not assume they are already adequately served.
  \item \textbf{Do not rely on settlement-count coverage statistics alone.} A population-weighted reading is essential given how large the divergence can be between the two (Section~\ref{sec:frame}); reports and dashboards built on this data should default to population-weighted figures, as this report does.
  \item \textbf{Treat the no-time-estimate population as a data-collection and network-improvement priority}, not as a resolved ``low access'' or ``adequate access'' category -- closing this measurement gap, e.g., through better rural road mapping or a wider sampling frame, is itself an actionable next step.
\end{enumerate}

Two methodological extensions would materially improve this analysis without changing its scope: a genuinely national (not per-department) road graph, to remove the cross-department limitation entirely; and a bulk, machine-readable official population source at the individual centro-poblado level, to reduce reliance on a district-level calibration target. Neither is required to trust the findings reported here, which have been validated at multiple levels (department, province, district, and total), but both would tighten the estimates further.
